\documentclass[12pt]{article}

\usepackage[margin=1in]{geometry}
\usepackage{booktabs}
\usepackage{array}
\usepackage{tabularx}
\usepackage{hyperref}
\usepackage{amsmath}

\title{Empirical Counterparts for MVPF Terms}
\author{Taco Prins, Lennard Schlattmann, and Daniel Jonas Schmidt}
\date{\today}

\begin{document}

\maketitle

\section*{Empirical Estimates by Term}

\begin{table}[h]
\centering
\small
\renewcommand{\arraystretch}{1}
\begin{tabularx}{\linewidth}{p{4.6cm}p{4.3cm}X}
\toprule
\textbf{Term} & \textbf{Estimate} & \textbf{Source} \\
\midrule
$I$ (take-up rate)
    & $\approx 30\%$ in high-risk areas
    & \href{https://www.fema.gov/openfema-data-page/nfip-residential-penetration-rates-v1}{FEMA OpenFEMA NFIP Penetration Rates}; \href{https://www.rand.org/content/dam/rand/pubs/technical_reports/2006/RAND_TR300.pdf}{Dixon et al.\ (2006)} \\
$\varepsilon$ (price semi-elasticity)
    & $-0.32$ (preferred); range $-0.17$ to $-1.38$
    & \href{https://faculty.wharton.upenn.edu/wp-content/uploads/2017/07/MismeasuringRisk_Mulder2021.pdf}{Mulder (2021)};
      \href{https://journalofcrr.com/research/03-07-gourevitch-et-al/}{Gourevitch et al.\ (2023)};
      \href{https://link.springer.com/article/10.1023/A:1007823631497}{Browne \& Hoyt (2000)};
      Gruber \& Solomon (2026) \\
$\partial I / \partial a$ (crowd-out)
    & $-3\%$ of coverage
    & \href{https://www.sciencedirect.com/science/article/abs/pii/S0095069617303479}{Kousky, Michel-Kerjan \& Raschky (2018)} \\
$F(q)$ (belief distribution)
    & Substantial heterogeneity
    & \href{https://www.nber.org/system/files/working_papers/w23854/w23854.pdf}{Bakkensen \& Barrage (2017)};
      \href{https://www.aeaweb.org/articles?id=10.1257/app.6.3.206}{Gallagher (2014)};
      \href{https://faculty.wharton.upenn.edu/wp-content/uploads/2017/07/MismeasuringRisk_Mulder2021.pdf}{Mulder (2021)} \\
$p$ (true flood probability)
    & $1\%$--$10\%$ per year
    & \href{https://zenodo.org/records/6459076}{First Street Foundation (Zenodo v2.0)};
      \href{https://www.fema.gov/flood-maps/national-flood-hazard-layer}{FEMA National Flood Hazard Layer} \\
$d$ (flood damage, conditional on flood)
    & \$44k--\$83k (avg.\ claim)
    & \href{https://www.fema.gov/openfema-data-page/fima-nfip-redacted-claims-v2}{OpenFEMA NFIP Redacted Claims};
      \href{https://www.floodsmart.gov/historical-nfip-claims-information-and-trends}{NFIP Historical Claims (FloodSmart)} \\
$d/w$ (damage-to-wealth ratio)
    & $10\%$--$16\%$
    & \href{https://www.fema.gov/openfema-data-page/fima-nfip-redacted-claims-v2}{OpenFEMA NFIP Redacted Claims};
      \href{https://www.floodsmart.gov/historical-nfip-claims-information-and-trends}{NFIP Historical Claims (FloodSmart)} \\
$w$ (household wealth)
    & \$193k (median); \$300k--\$500k (home value)
    & \href{https://fortune.com/2023/10/18/us-household-wealth-increase-survey-of-consumer-finances-federal-reserve/}{Federal Reserve SCF (2022)};
      \href{https://www.census.gov/programs-surveys/ahs.html}{American Housing Survey} \\
$s$ (current subsidy rate)
    & $\approx 47\%$ (median policyholder)
    & \href{https://www.gao.gov/assets/gao-23-105977.pdf}{GAO (2023)};
      \href{https://www.fema.gov/flood-insurance/risk-rating}{FEMA Risk Rating 2.0} \\
$a$ (current relief fraction)
    & $\approx 14\%$ (preferred); range $4\%$--$7\%$
    & Gruber \& Solomon (2026): avg.\ IHP payout \$20k vs.\ avg.\ damage \$144k in flood-damage-weighted sample;
      lower range from avg.\ FEMA IA grant (\$3k) over avg.\ claim (\$44k--\$83k);
      \href{https://www.sciencedirect.com/science/article/abs/pii/S0095069617303479}{Kousky et al.\ (2018)} \\
$\lambda$ (shadow value of budget)
    & $0.5$--$2$
    & \href{https://academic.oup.com/qje/article/135/3/1209/5781614}{Hendren \& Sprung-Keyser (2020)} \\
\bottomrule
\end{tabularx}
\caption{Empirical counterparts for the terms in MVPF$_s$ and MVPF$_a$. All estimates are from US data except where noted. Preferred values for $\varepsilon$ and $a$ are from Gruber \& Solomon (2026). The crowd-out estimate from Kousky et al.\ (2018) is identified on the intensive margin (coverage amount) rather than the binary take-up margin. The model calibration of $d/w = 0.15$ is consistent with average NFIP claims (\$44,000--\$83,000) relative to average US home values (\$300,000--\$500,000).}
\end{table}

\section*{Notes on Identification}

\paragraph{Price semi-elasticity $\varepsilon$.}
We define $\varepsilon \equiv \partial I/\partial\ln\pi < 0$ as the semi-elasticity of take-up with respect to the log-premium. Since $\ln\pi = \ln((1-s)pd)$, we have $\partial\ln\pi/\partial s = -1/(1-s)$ and therefore
\[
    \frac{\partial I}{\partial s} = -\frac{\varepsilon}{1-s}.
\]
Mulder (2021) exploits variation from outdated FEMA flood maps that misclassify high-risk homes as low-risk, yielding $|\varepsilon| \approx 0.17$.
Gourevitch et al.\ (2023) exploit the NFIP's Risk Rating~2.0 reform and find a wider range of $|\varepsilon|$ from 0.15 to 1.38 across subgroups.
Gruber \& Solomon (2026) estimate $\varepsilon = -0.32$ using within-property variation in billed premiums as policies transition toward RR2.0 full-risk pricing; they explicitly report this as a semi-elasticity.

\paragraph{Crowd-out $\partial I / \partial a$.}
Direct quasi-experimental estimates of $\partial I/\partial a$ are rare because variation in disaster relief generosity is uncommon.
Kousky et al.\ (2018) link FEMA Individual Assistance grants (2000--2011) to NFIP policy data and find that receiving a grant reduces coverage the following year by \$4{,}000--\$5{,}000 ($\approx 3\%$ of mean coverage).
Alternatively, using the ratio trick (eq.~7 of the model notes),
\[
    \frac{\partial I/\partial a}{\partial I/\partial s} = -\frac{q^*}{p}\cdot\frac{u'(c_{U,F})}{u'(c_I)},
\]
$\partial I/\partial a$ can be recovered from the estimated $\varepsilon$ together with structural utility assumptions.

\paragraph{Belief distribution $F(q)$.}
Gallagher (2014) documents that take-up spikes after a flood and then decays, consistent with $q < p$ on average and with beliefs that revert toward prior underestimation over time.
Bakkensen and Barrage (2017) document large cross-sectional heterogeneity in subjective flood probabilities among US coastal homeowners.
Mulder (2021) measures subjective flood risk at the property level alongside objective risk, enabling direct estimation of the belief gap $p - q$.

\paragraph{Risk aversion $\gamma$.}
Chetty (2006) identifies $\gamma \approx 1$ from labor supply responses, providing a lower bound.
The broader empirical literature yields estimates in $[1, 4]$; we use $\gamma = 2$ which is in the middle of this range.

\paragraph{Shadow value $\lambda$.}
At the interior optimum, MVPF$_s$ = MVPF$_a$ = $\lambda$.
Hendren and Sprung-Keyser (2020) compute MVPFs for 133 US policy changes; their estimates for adult-targeted social insurance programs cluster between 0.5 and 2, providing a plausible external benchmark for $\lambda$.

\section*{Numerical MVPF Estimates at Current US Policy}

Using the preferred empirical parameter values above, we evaluate both MVPFs at the current US policy mix ($s = 0.47$, $a = 0.14$, $d/w = 0.15$, $\varepsilon = -0.32$).
The crowd-out response $\partial I/\partial a$ is recovered via the ratio trick.
Table~\ref{tab:mvpf} reports all intermediate quantities and the resulting MVPFs.

\begin{table}[h]
\centering
\renewcommand{\arraystretch}{1.4}
\begin{tabular}{lc}
\toprule
& \textbf{Value} \\
\midrule
\multicolumn{2}{l}{\textit{Parameters}} \\
$p$           & $0.02$ \\
$d/w$         & $0.15$ \\
$s$           & $0.47$ \\
$a$           & $0.14$ \\
$I$           & $0.30$ \\
$\varepsilon$ & $-0.32$ \\
$\gamma$      & $2.0$ \\
\midrule
\multicolumn{2}{l}{\textit{Consumption levels}} \\
$c_I$     & $0.9984$ \\
$c_{U,F}$ & $0.8710$ \\
\midrule
\multicolumn{2}{l}{\textit{Intermediate quantities}} \\
$q^*$                         & $0.0108$ \\
$(p - q^*)\,\Delta u$         & $0.00137$ \\
$u'(c_I)$                     & $1.0032$ \\
$u'(c_{U,F})$                 & $1.3181$ \\
$\partial I/\partial s$       & $0.604$ \\
$\partial I/\partial a$       & $-0.427$ \\
\midrule
\multicolumn{2}{l}{\textit{MVPFs}} \\
$\mathrm{MVPF}_s$ & $\mathbf{1.16}$ \\
$\mathrm{MVPF}_a$ & $\mathbf{1.30}$ \\
\bottomrule
\end{tabular}
\caption{MVPF estimates at current US policy ($s = 0.47$, $a = 0.14$, $I = 0.30$, $\varepsilon = -0.32$, $\gamma = 2$, $p = 0.02$, $d/w = 0.15$). Preferred parameter values updated to reflect Gruber \& Solomon (2026): $\varepsilon = -0.32$ replaces the earlier $-0.17$ from Mulder (2021); $a = 0.14$ replaces $0.055$, based on the flood-damage-weighted IHP payout of \$20k against average damage of \$144k. The consumption level $c_{U,F} = w - (1-a)d = 1 - 0.86 \times 0.15 = 0.8710$. The semi-elasticity $\varepsilon = -0.32$ enters via $\partial I/\partial s = -\varepsilon/(1-s) = 0.604$; the crowd-out $\partial I/\partial a = -0.427$ is recovered from the ratio trick.}
\label{tab:mvpf}
\end{table}

With the updated $a = 0.14$, flood victims retain more consumption ($c_{U,F} = 0.871$ vs.\ $0.858$ previously), reducing their marginal utility advantage over insured households. Interpreting $\varepsilon = -0.32$ as the semi-elasticity $\partial I/\partial\ln\pi$ yields $\partial I/\partial s = -\varepsilon/(1-s) = 0.604$, far larger than the old formula (which incorrectly treated $\varepsilon$ as a true elasticity $(1/I)(\partial I/\partial\pi)$) implied. Both MVPFs exceed 1: $\mathrm{MVPF}_s = 1.16$ and $\mathrm{MVPF}_a = 1.30$, indicating that both instruments generate positive net welfare gains at the current policy mix, with disaster relief having a somewhat higher return.

As a cross-check on $a = 0.14$: with $p = 0.02$, $d/w = 0.15$, and $w = \$500{,}000$, the model implies a fiscal spillover of $a \times p \times d = 0.14 \times 0.02 \times \$75{,}000 = \$210$ per house per percentage-point increase in coverage. This is almost exactly Gruber \& Solomon's IV estimate of \$175--\$203, providing an external validation of the updated $a$.

\paragraph{Utility function.}
All computations use CRRA utility with $\gamma = 2$:
\[
    u(c) = \frac{c^{1-\gamma}}{1-\gamma} = -\frac{1}{c}, \qquad u'(c) = c^{-\gamma} = \frac{1}{c^2}.
\]
Utility is defined over total final consumption, with wealth normalized to $w = 1$.
The consumption levels entering the MVPF formulas are therefore
\begin{align*}
    c_I     &= w - (1-s)\,pd = 1 - 0.53 \times 0.003 = 0.9984, \\
    c_{U,F} &= w - (1-a)\,d = 1 - 0.86 \times 0.15 = 0.8710.
\end{align*}
Because the actuarially fair premium $pd = 0.02 \times 0.15 = 0.003$ is small, insured households barely move away from $w = 1$, so $u'(c_I) \approx 1$.
This means the inframarginal welfare contribution to $\mathrm{MVPF}_s$ is approximately 1; the numerator exceeds 1 because the belief-gap internality $(p-q^*)f(q^*)$ adds to the effective recipient count.
The gap $\mathrm{MVPF}_a - \mathrm{MVPF}_s$ is driven by the higher marginal utility of flood victims, though it is smaller under $a = 0.14$ than under the old $a = 0.055$ because better relief leaves flood victims closer to full consumption.

\end{document}
